矩阵用两个索引组织信息：行与列、输入与输出、样本与特征。但彩色视频同时有时间、高度、宽度和颜色，推荐数据可能同时依赖用户、商品与场景；若把它们粗暴摊平成矩阵，计算虽然熟悉，``哪些方向原本属于同一种变化''却被藏进了人为编号。

张量把矩阵的双索引结构推广到多个方向，但``多维数组''只是选定基后的坐标外观。真正不依赖坐标的对象是多线性映射或张量积中的元素。这个区分决定了换基时分量怎样联动，也解释了张量分解为什么不是把矩阵算法机械加一个索引。

\subsection*{一个两分钟试算}

取最简单的 \(2\times 2\times 2\) 三阶张量，把它看成两个 \(2\times 2\) 的前片：
\[
X_{:,:,1}=
\begin{bmatrix}1&2\\3&4\end{bmatrix},
\qquad
X_{:,:,2}=
\begin{bmatrix}5&6\\7&8\end{bmatrix}.
\]
这是一个多维数组，记下 \(2\times 2\times 2=8\) 个数。按第一模``摊平''成 \(2\times 4\) 矩阵：
\[
X_{(1)}=
\begin{bmatrix}
1&2&5&6\\
3&4&7&8
\end{bmatrix},
\]
秩为 2。按第二模摊平：
\[
X_{(2)}=
\begin{bmatrix}
1&3&5&7\\
2&4&6&8
\end{bmatrix},
\]
秩也为 2。再按第三模摊平：
\[
X_{(3)}=
\begin{bmatrix}
1&3&2&4\\
5&7&6&8
\end{bmatrix},
\]
秩仍是 2。三个展开矩阵恰好给出了同一个多线性秩 \((2,2,2)\)——但这是凑巧。一个 \(2\times 3\times 5\) 的张量，三个展开矩阵的秩完全可以分别是 2、3、4。CP 秩（最少秩一张量项数）又可能完全不同。这就是为什么不能简单把矩阵秩搬到高维：同一个数组，沿不同模摊平，看到的结构不同。

\medskip
\begin{enumerate}
\tightlist
\item
  张量是多线性映射，也是多维数组：从双线性型的矩阵表示出发，引出阶、维度、外积、缩并与换基，分清对象本身和它的坐标数组。
\item
  张量分解初步：CP、Tucker 与 SVD 的关系：把秩一外积推广到多路数据，用 CP 寻找少数可分离成分，用 Tucker 寻找各模子空间与核心交互，并说明高阶秩不再拥有矩阵秩的全部好性质。
\end{enumerate}

本单元只讨论有限维实向量空间和基础分解。指标升降、微分几何中的张量场、量子多体网络与大规模优化各有额外结构；它们共享多线性语言，却不能只凭数组形状互相替代。

\subsection{怎么读这一章}

第一次阅读可只抓四个动作：为每条轴命名，用外积生成可分离结构，用缩并消去配对索引，用模乘法只改变指定轴的坐标。掌握这四件事后，就可直接读第二篇中 CP 与 Tucker 的对比：前者假设少数跨模可分离成分，后者假设每个模各有低维子空间，再用核心记录它们的交互。

若只想为深度学习回补 shape 直觉，可先读第一篇的``阶与维度''、``模乘法''和``缩并''；但要记住，软件框架称为 tensor 的对象首先是带形状的数组，只有当每条轴对应的空间与换基规则也被声明时，它才表示数学上不依赖坐标的张量。

本单元的第一层停点是：能说清``阶''、``各模维度''、``CP 秩''与``多线性秩''为何是不同对象，并能指出矩阵 SVD 的正交、最佳截断与唯一奇异值序列为何不是高阶张量的默认性质。

\subsection{本章篇目}

\begin{itemize}
\tightlist
\item \hyperref[sec:03-Linear-Algebra/09-Tensors/01]{``张量是多线性映射，也是多维数组''}；
\item \hyperref[sec:03-Linear-Algebra/09-Tensors/02]{``张量分解初步：CP、Tucker 与 SVD 的关系''}。
\end{itemize}
